
\documentclass[]{standalone}
\usepackage[T1]{fontenc}

\usepackage{tikz}
\usepackage{amsmath, amsfonts}
\usetikzlibrary{arrows, shapes.geometric, calc, positioning, decorations.pathmorphing}

\makeatletter
\def\input@path{{./includes/}{./includes/ut/diags/}}
\makeatother

% start -> interop? -> PoA works? -> PoS needed

% multiline helper
\NewDocumentCommand{\mline}{m}{%
\begin{tabular}{c}#1\end{tabular}%
}

\begin{document}
\tikzstyle{node} = [rectangle, rounded corners, minimum width=2cm, minimum height=0.8cm, text centered, draw=black]
\tikzstyle{arrow} = [-stealth, shorten >= 2pt, shorten <= 2pt]
\tikzstyle{conflict} = [stealth-stealth, decorate, decoration={zigzag,pre=lineto,pre length=2mm,post=lineto,post length=1mm,segment length=5mm,amplitude=1mm}]
\begin{tikzpicture}[node distance=1.25cm]
    \node (start) [node] {\mline{Need a\\Dapp-chain?}};
    \node (interop) [node, right = of start, xshift=0mm] {\mline{Requires\\interop. with\\Simplex?}};
    \node (poa) [node, right = of interop, xshift=0mm] {\mline{Does PoA\\work for\\use case?}};
    % \node (pos) [right = of poa, xshift=0mm] {}; % {\mline{Requires\\PoS}};

    \node (usePoa) [node, below = of poa] {\mline{Use PoA}};
    \node (noInterop) [node] at (interop |- usePoa) {\mline{Whatever}};
    \node (usePos) [node, right = of usePoa] {\mline{Use PoS}};
    \node (trCorner) [] at (poa -| usePos) {};

    \draw [arrow] (start) -- (interop);
    \draw [arrow] (interop) -- node (yes1) [above] {Yes} (poa);
    \draw [arrow] (poa) -- node (no2) [above] {No} (poa -| usePos) -- (usePos);
    \draw [arrow] (interop) -- node (no1) [right] {No} (noInterop);
    \draw [arrow] (poa) -- node (yes2) [right] {Yes} (usePoa);
    \draw [arrow] (start) -- node (yes0) [above] {Yes} (interop);
    \draw [arrow] (start) -- node (no0) [right] {No} (start |- noInterop) -- (noInterop);


    % position to right of midpoint between start and mine
    % \node (restart) [right = of sm-mid.east, xshift=8mm] {};
    % \node (restartLabel) [right = of restart.west, xshift=-12mm] {\begin{tabular}{c} Restart or\\ Update \end{tabular}};
    % \node (restartLabel) [node, right = of sm-mid.east, xshift=8mm] {\begin{tabular}{c} Restart or Update \end{tabular}};

    % \draw [arrow] (mine) -- (mine -| restart.west) -- node (restartLabel) [right, xshift=-2mm] {\begin{tabular}{c} Restart or\\ Update \end{tabular}} (restart.west |- start) -- (start);

    % \node (broadcast) [node, below = of mine] {Broadcast the block};
    % \draw [arrow] (mine) -- node (mb-lab) [pos=0.5, right] {Find PoW} (broadcast);


    % \node (goal) [node] {\begin{tabular}{c} Safely increase capacity \\ via merged mining \end{tabular}};

    % \node (m1) [right = of goal.east] {};
    % \node (mr1) [node, above = of goal.east -| m1] {Stay decentralized};
    % \node (mr2) [node, below = of goal.east -| m1] {Securely add more chains};
    % \draw [arrow] (goal) -- (mr1);
    % \draw [arrow] (goal) -- (mr2);

    % \node (arp1) [node, right = of mr1.east, xshift=2mm] {\begin{tabular}{c} Users don't have to \\ validate all chains \end{tabular}};
    % \node (arp2) [node] at (mr2 -| arp1) {\begin{tabular}{c} Miners must \\ validate all chains \end{tabular}};
    % \draw [arrow] (mr1) -- (arp1);
    % \draw [arrow] (mr2) -- (arp2);

    % \draw [conflict] (arp1) -- node (c) [anchor=west, inner sep=5pt] {\begin{tabular}{c} \textbf{Conflict} \\ Miners are Users \end{tabular}} (arp2);


    %\node (conflict) [node] at ([xshift=1cm]c.west) {\textbf{Conflict}};
\end{tikzpicture}
\end{document}
